% Algebra 2 Lesson 4-5 fast-track packet: Solving Rational Equations
% (enVision A2 2024 TE, Example 2 and Example 4 only; fast-lesson draft 2026-10-31, re-cut toward the Topic 4 test; PLC has not met on this lesson)
% Topic 4 test (a2_T4_assessment_TE.tex): items 7 and 13 work rate (Example 2, kept); items 9 and 11 solve and check (Example 4, kept);
% items 6 and 17 product denominator x(x-2) or x(x-3) (item 6's root is extraneous, no solution; item 17's root x = 2 is valid) (Example 3, cut; rehearsed by TE Practice 24, SEE EXAMPLE 4, in the Practice block).
% Quiz map (TE "4-5 Lesson Quiz", all DOK 2, skill A87): item 2 work rate (Example 2; graded); item 4 extraneous solution
% (Example 4, Part B is its structural twin; graded); item 3 a variable denominator on each side (Example 1, cut; proposed bonus,
% rehearsed by TE Practice 32 and Practice 24 in the Practice block); items 1 (Example 1) and 5 (Example 5) full Savvas quiz only.
% Shared preamble for the "fast" (hacked-down) Algebra 2 lesson packets, Oct 2026.
% Each packet: % Shared preamble for the "fast" (hacked-down) Algebra 2 lesson packets, Oct 2026.
% Each packet: % Shared preamble for the "fast" (hacked-down) Algebra 2 lesson packets, Oct 2026.
% Each packet: \input{fast-packet-template} then \begin{document} ... \end{document}.
\documentclass[12pt]{article}
\usepackage[margin=0.7in]{geometry}
\usepackage{amsmath,amssymb}
\usepackage{booktabs,tabularx,array}
\usepackage{enumitem}
\usepackage[T1]{fontenc}
\usepackage[utf8]{inputenc}
\usepackage{lmodern}
\usepackage{tikz}
\usepackage{pgfplots}
\pgfplotsset{compat=1.18}
\usepackage{xcolor}
\setlist{nosep,leftmargin=1.4em}
\setlength{\parindent}{0pt}
\setlength{\parskip}{6pt}
\renewcommand{\arraystretch}{1.35}
\newcolumntype{Y}{>{\raggedright\arraybackslash}X}
\pagestyle{empty}

% blanks and workspace
\newcommand{\blank}[1][2.6cm]{\rule[-2pt]{#1}{0.5pt}}
\newcommand{\lines}[1][1]{\par\vspace{2pt}\foreach \i in {1,...,#1}{\rule{0pt}{22pt}\rule{\linewidth}{0.4pt}\par}}
\newcommand{\wspart}[2]{\vspace{4pt}{\large\bfseries #1: #2}\par\vspace{-6pt}\rule{\textwidth}{0.4pt}\par}

% header: \packethead{Lesson 1-3}{Piecewise-Defined Functions}
\newcommand{\packethead}[2]{%
{\LARGE\bfseries #2}\hfill Name \blank[5cm]\\[2pt]
Algebra 2 \quad #1 (fast track) \hfill Date \blank[2.5cm] \quad Period \blank[1.2cm]\par\vspace{2pt}}

% boxed goals: \goals{learning goal}{language goal}
\newcommand{\goals}[2]{%
\fbox{\parbox{\dimexpr\textwidth-2\fboxsep-2\fboxrule}{%
\textbf{Learning goal:} #1\\[2pt]
\textbf{Language goal:} #2}}\par}

% quiz target box: \quiztarget{text}
\newcommand{\quiztarget}[1]{%
\vspace{4pt}\colorbox{gray!12}{\parbox{\dimexpr\textwidth-2\fboxsep}{\textbf{On the quiz:} #1}}\par}

% IXL table row helper: code & skill & which quiz question it feeds
\newcommand{\ixlrow}[3]{\texttt{#1} & #2 & #3 \\}

% one grid style for every graph; override xmin/xmax/... per plot
\pgfplotsset{sheetgrid/.style={
  axis lines=middle, grid=both, grid style={gray!45},
  xmin=-10.5, xmax=10.5, ymin=-10.5, ymax=10.5,
  xtick={-10,-8,...,10}, ytick={-10,-8,...,10},
  x=0.28cm, y=0.28cm,
  ticklabel style={font=\scriptsize, fill=white, inner sep=1pt},
  xlabel={$x$}, ylabel={$y$},
  axis line style={thick}, clip=true}}
% unit grid for small windows
\pgfplotsset{unitgrid/.style={sheetgrid,
  xmin=-5.5, xmax=5.5, ymin=-5.5, ymax=5.5,
  xtick={-5,...,5}, ytick={-5,...,5}, x=0.5cm, y=0.5cm}}
% open and closed endpoints
\tikzset{closed/.style={circle, fill, inner sep=1.6pt},
         open/.style={circle, draw, fill=white, inner sep=1.4pt}}
 then \begin{document} ... \end{document}.
\documentclass[12pt]{article}
\usepackage[margin=0.7in]{geometry}
\usepackage{amsmath,amssymb}
\usepackage{booktabs,tabularx,array}
\usepackage{enumitem}
\usepackage[T1]{fontenc}
\usepackage[utf8]{inputenc}
\usepackage{lmodern}
\usepackage{tikz}
\usepackage{pgfplots}
\pgfplotsset{compat=1.18}
\usepackage{xcolor}
\setlist{nosep,leftmargin=1.4em}
\setlength{\parindent}{0pt}
\setlength{\parskip}{6pt}
\renewcommand{\arraystretch}{1.35}
\newcolumntype{Y}{>{\raggedright\arraybackslash}X}
\pagestyle{empty}

% blanks and workspace
\newcommand{\blank}[1][2.6cm]{\rule[-2pt]{#1}{0.5pt}}
\newcommand{\lines}[1][1]{\par\vspace{2pt}\foreach \i in {1,...,#1}{\rule{0pt}{22pt}\rule{\linewidth}{0.4pt}\par}}
\newcommand{\wspart}[2]{\vspace{4pt}{\large\bfseries #1: #2}\par\vspace{-6pt}\rule{\textwidth}{0.4pt}\par}

% header: \packethead{Lesson 1-3}{Piecewise-Defined Functions}
\newcommand{\packethead}[2]{%
{\LARGE\bfseries #2}\hfill Name \blank[5cm]\\[2pt]
Algebra 2 \quad #1 (fast track) \hfill Date \blank[2.5cm] \quad Period \blank[1.2cm]\par\vspace{2pt}}

% boxed goals: \goals{learning goal}{language goal}
\newcommand{\goals}[2]{%
\fbox{\parbox{\dimexpr\textwidth-2\fboxsep-2\fboxrule}{%
\textbf{Learning goal:} #1\\[2pt]
\textbf{Language goal:} #2}}\par}

% quiz target box: \quiztarget{text}
\newcommand{\quiztarget}[1]{%
\vspace{4pt}\colorbox{gray!12}{\parbox{\dimexpr\textwidth-2\fboxsep}{\textbf{On the quiz:} #1}}\par}

% IXL table row helper: code & skill & which quiz question it feeds
\newcommand{\ixlrow}[3]{\texttt{#1} & #2 & #3 \\}

% one grid style for every graph; override xmin/xmax/... per plot
\pgfplotsset{sheetgrid/.style={
  axis lines=middle, grid=both, grid style={gray!45},
  xmin=-10.5, xmax=10.5, ymin=-10.5, ymax=10.5,
  xtick={-10,-8,...,10}, ytick={-10,-8,...,10},
  x=0.28cm, y=0.28cm,
  ticklabel style={font=\scriptsize, fill=white, inner sep=1pt},
  xlabel={$x$}, ylabel={$y$},
  axis line style={thick}, clip=true}}
% unit grid for small windows
\pgfplotsset{unitgrid/.style={sheetgrid,
  xmin=-5.5, xmax=5.5, ymin=-5.5, ymax=5.5,
  xtick={-5,...,5}, ytick={-5,...,5}, x=0.5cm, y=0.5cm}}
% open and closed endpoints
\tikzset{closed/.style={circle, fill, inner sep=1.6pt},
         open/.style={circle, draw, fill=white, inner sep=1.4pt}}
 then \begin{document} ... \end{document}.
\documentclass[12pt]{article}
\usepackage[margin=0.7in]{geometry}
\usepackage{amsmath,amssymb}
\usepackage{booktabs,tabularx,array}
\usepackage{enumitem}
\usepackage[T1]{fontenc}
\usepackage[utf8]{inputenc}
\usepackage{lmodern}
\usepackage{tikz}
\usepackage{pgfplots}
\pgfplotsset{compat=1.18}
\usepackage{xcolor}
\setlist{nosep,leftmargin=1.4em}
\setlength{\parindent}{0pt}
\setlength{\parskip}{6pt}
\renewcommand{\arraystretch}{1.35}
\newcolumntype{Y}{>{\raggedright\arraybackslash}X}
\pagestyle{empty}

% blanks and workspace
\newcommand{\blank}[1][2.6cm]{\rule[-2pt]{#1}{0.5pt}}
\newcommand{\lines}[1][1]{\par\vspace{2pt}\foreach \i in {1,...,#1}{\rule{0pt}{22pt}\rule{\linewidth}{0.4pt}\par}}
\newcommand{\wspart}[2]{\vspace{4pt}{\large\bfseries #1: #2}\par\vspace{-6pt}\rule{\textwidth}{0.4pt}\par}

% header: \packethead{Lesson 1-3}{Piecewise-Defined Functions}
\newcommand{\packethead}[2]{%
{\LARGE\bfseries #2}\hfill Name \blank[5cm]\\[2pt]
Algebra 2 \quad #1 (fast track) \hfill Date \blank[2.5cm] \quad Period \blank[1.2cm]\par\vspace{2pt}}

% boxed goals: \goals{learning goal}{language goal}
\newcommand{\goals}[2]{%
\fbox{\parbox{\dimexpr\textwidth-2\fboxsep-2\fboxrule}{%
\textbf{Learning goal:} #1\\[2pt]
\textbf{Language goal:} #2}}\par}

% quiz target box: \quiztarget{text}
\newcommand{\quiztarget}[1]{%
\vspace{4pt}\colorbox{gray!12}{\parbox{\dimexpr\textwidth-2\fboxsep}{\textbf{On the quiz:} #1}}\par}

% IXL table row helper: code & skill & which quiz question it feeds
\newcommand{\ixlrow}[3]{\texttt{#1} & #2 & #3 \\}

% one grid style for every graph; override xmin/xmax/... per plot
\pgfplotsset{sheetgrid/.style={
  axis lines=middle, grid=both, grid style={gray!45},
  xmin=-10.5, xmax=10.5, ymin=-10.5, ymax=10.5,
  xtick={-10,-8,...,10}, ytick={-10,-8,...,10},
  x=0.28cm, y=0.28cm,
  ticklabel style={font=\scriptsize, fill=white, inner sep=1pt},
  xlabel={$x$}, ylabel={$y$},
  axis line style={thick}, clip=true}}
% unit grid for small windows
\pgfplotsset{unitgrid/.style={sheetgrid,
  xmin=-5.5, xmax=5.5, ymin=-5.5, ymax=5.5,
  xtick={-5,...,5}, ytick={-5,...,5}, x=0.5cm, y=0.5cm}}
% open and closed endpoints
\tikzset{closed/.style={circle, fill, inner sep=1.6pt},
         open/.style={circle, draw, fill=white, inner sep=1.4pt}}

\begin{document}

\packethead{Lesson 4-5}{Solving Rational Equations}
\goals{I can solve a rational equation, including a work-rate problem, by multiplying both sides by a common denominator, and I can identify any extraneous solutions.}{I can explain why a solution is or is not valid using the words rational equation, common denominator, rate, domain and extraneous solution.}
\vspace{-4pt}{\footnotesize Spanish: \emph{ecuación racional} (rational equation), \emph{solución extraña} (extraneous solution), \emph{dominio} (domain).}

\vspace{-4pt}\quiztarget{(2) graded: add the workers' \textbf{rates} (Example 2, Practice 1). (4) graded: clear the fractions, solve, cross out the \textbf{extraneous solution} (Example 4). (3) bonus: a variable denominator on each side (Practice 2, 3). (1), (5): full Savvas quiz only.}
\vspace{-4pt}\topictarget{4}{Questions 7 and 13 ask for a time alone or together (Example 2). Questions 9 and 11 ask you to solve and check (Example 4). Questions 6 and 17 (a denominator like $x(x-3)$) come from Example 3, which this packet skips; Practice 2 rehearses them.}

\wspart{Example 2}{Solve a Work-Rate Problem}
% TE Example 2 statement, lightly compressed; "Cheyenne works twice as fast as Arthur" is the TE illustration's callout.
Arthur and Cheyenne are making Chinese paper cutting designs to give as gifts for the coming year. Together they can complete the job in 24 hours. Cheyenne works twice as fast as Arthur. How long would it take Cheyenne to complete the job if she was working alone?

\begin{minipage}[t]{0.50\textwidth}
\vspace{0pt}\raggedright
\textbf{Step 1} Determine the work-rates. The job done per hour is the \textbf{rate}. Let $x$ = Arthur's hours alone.\\[2pt]
Arthur: $\frac{1\text{ job}}{x\text{ hours}}=\frac1x$ job per hour.\\[2pt]
Cheyenne: $\frac{1\text{ job}}{\frac12x\text{ hours}}=\frac2x$ job per hour.\\[2pt]
Together: $\frac{1\text{ job}}{24\text{ hours}}=\frac1{24}$ job per hour.\\[3pt]
% TE Example 2 callout (COMMON ERROR), compressed.
{\small\textbf{Watch out:} twice as fast means \emph{half} the time, $\tfrac12x$, not $2x$.}
\end{minipage}\hfill
\begin{minipage}[t]{0.46\textwidth}
\vspace{0pt}\raggedright
\textbf{Step 2} Write the \textbf{rational equation} for their rates together. Solve. {\small Domain: $x\neq0$.}
\vspace{-6pt}
\begin{align*}
\frac1x+\frac2x&=\frac1{24}\\
24x\left(\frac1x+\frac2x\right)&=\left(\frac1{24}\right)24x\\
24+48&=x \quad\Rightarrow\quad x=72
\end{align*}
\textbf{Check:} $\tfrac1{72}+\tfrac2{72}=\tfrac3{72}=\tfrac1{24}$ \checkmark\\[2pt]
\fbox{Arthur: 72 hours. Cheyenne alone: 36 hours.}
\end{minipage}

\vspace{3pt}
\textbf{Say it:} Arthur's \blank[1.4cm] is $\tfrac1x$ job per hour. I multiplied both sides of the \blank[3cm] by $24x$, the \blank[3.4cm]. The answer 72 is in the \blank[1.6cm].

\wspart{Try It 2}{Solve a work-rate problem}
It takes 12 hours to fill a pool with two pipes, where the water in one pipe flows three times as fast as the other pipe. How long will it take the slower pipe to fill the pool by itself?

\textbf{Hint:} $x$ = slower pipe's hours alone. \ Rates: slower \blank[1.2cm], faster \blank[1.2cm]. \ Equation: \blank[2.4cm]
\lines[2]
The slower pipe alone: \blank[2cm] hours

\newpage
\wspart{Example 4}{Check for Extraneous Solutions}
What are the solutions to the following equations? First find the \textbf{domain}: the values that make a denominator 0 can never be solutions. A value that the cleared equation gives but the domain leaves out is an \textbf{extraneous solution}.

\begin{minipage}[t]{0.45\textwidth}
\vspace{0pt}\raggedright
\textbf{A.} $\dfrac{5x}{x-2}=7+\dfrac{10}{x-2}$

{\small Domain: $x\neq2$.\\ Common denominator: $x-2$.}
\begin{align*}
(x-2)\left(\frac{5x}{x-2}\right)&=\left(7+\frac{10}{x-2}\right)(x-2)\\
5x&=7(x-2)+10\\
5x&=7x-14+10\\
-2x&=-4\\
x&=2
\end{align*}
\textbf{Check:} $x=2$ makes the denominator $x-2$ equal to 0, so 2 is an \textbf{extraneous solution}.

\fbox{This equation has \textbf{no solution}.}
\end{minipage}\hfill
\begin{minipage}[t]{0.51\textwidth}
\vspace{0pt}\raggedright
\textbf{B.} $\dfrac{3}{x-3}=\dfrac{x}{x-3}-\dfrac{x}{4}$

{\small Domain: $x\neq3$. \ Common denominator: $4(x-3)$.}
\begin{align*}
4(x-3)\left(\frac{3}{x-3}\right)&=\left(\frac{x}{x-3}-\frac{x}{4}\right)4(x-3)\\
4(3)&=4(x)-x(x-3)\\
12&=4x-x^2+3x\\
x^2-7x+12&=0 \quad \text{\small standard form}\\
(x-3)(x-4)&=0 \quad \text{\small factor}\\
x=3 \quad\text{or}\quad x&=4
\end{align*}
\textbf{Check:} $x=3$ makes $x-3=0$, so 3 is extraneous.\\
$x=4$: \ $\dfrac{3}{1}=3$ and $\dfrac{4}{1}-\dfrac{4}{4}=4-1=3$ \checkmark

\fbox{The only solution is $x=4$; \ 3 is extraneous.}
\end{minipage}

\vspace{2pt}

\textbf{Say it:} In Part B, $x=3$ is an \blank[3.4cm] because it makes the \blank[2.6cm] equal to 0, so it is not in the \blank[2cm] of the original equation. The only solution is \blank[1.6cm].

\wspart{Try It 4}{What are the solutions to the following equations?}
\begin{minipage}[t]{0.48\textwidth}
\vspace{0pt}
\textbf{a.} $x+\dfrac{6}{x-3}=\dfrac{2x}{x-3}$

Domain: $x\neq$\,\blank[1.2cm] \quad Multiply by: \blank[1.6cm]
\lines[4]
Solutions of the cleared equation: \blank[2.4cm]\\[4pt]
Extraneous: \blank[1.6cm] \quad Solution: \blank[1.6cm]
\end{minipage}\hfill
\begin{minipage}[t]{0.48\textwidth}
\vspace{0pt}
\textbf{b.} $\dfrac{x^2}{x+5}=\dfrac{25}{x+5}$

Domain: $x\neq$\,\blank[1.2cm] \quad Multiply by: \blank[1.6cm]
\lines[4]
Solutions of the cleared equation: \blank[2.4cm]\\[4pt]
Extraneous: \blank[1.6cm] \quad Solution: \blank[1.6cm]
\end{minipage}

\vspace{4pt}
In part b, the cleared equation has two solutions. Why is only one of them a solution of the original equation? Use the word \emph{extraneous}.
\lines[2]

% R5 Practice, part 1: rehearses graded Lesson Quiz item 2 (Aaron 6 h and Anna 3 h; how long together?): Example 2 and Try It 2
% find one worker's time from the time together, never the time together from two times alone. TE Practice & Problem Solving 26
% (Example 2, DOK 2), verbatim: same structure, different setting and numbers. It is also the shape of Topic 4 test item 13 (Faucets A and B).
% Part e is TE Practice 8 (UNDERSTAND: Reason, DOK 3, Example 2), verbatim.
% Part 2: rehearses Topic 4 test items 6 and 17 (a product denominator; item 6's root is extraneous, item 17's root x = 2 is valid; Example 3, cut). TE Practice 24 (SEE EXAMPLE 4, DOK 1), verbatim; TE answer x = -2.
% Part 3: rehearses bonus Lesson Quiz item 3 (3/(x-1) = 4/(2x+3)), a variable denominator on each side.
% TE Practice & Problem Solving 32 (SAT/ACT, SEE EXAMPLE 1, DOK 1), verbatim; TE answer D (x = 11).
\wspart{Practice}{Working together, then two denominators}
\textbf{1. Working together (Practice 26).} Kenji can finish a puzzle in 2 hours working alone. Oscar can finish the same puzzle in 3 hours working alone. How long would it take Oscar and Kenji to finish the puzzle if they worked on it together?

\textbf{a.} Each rate is 1 puzzle divided by the hours. \ Kenji: \blank[1.4cm] puzzle per hour. \ Oscar: \blank[1.4cm] puzzle per hour.

\textbf{b.} Let $t$ = hours working together. Their rates add up to $\tfrac1t$: \quad \blank[1.4cm] $+$ \blank[1.4cm] $=$ \blank[1.4cm]

\begin{minipage}[t]{0.48\textwidth}
\vspace{0pt}\raggedright
\textbf{c.} Multiply both sides by the common denominator \blank[1.4cm] and solve.
\lines[3]
\textbf{d.} $t=$\,\blank[2cm] hours
\end{minipage}\hfill
\begin{minipage}[t]{0.48\textwidth}
\vspace{0pt}
\textbf{e. (Practice 8)} If you solve a work-rate problem and your solution, which represents the amount of time it would take working together, exceeds the individual working alone times that are given, then how do you know your solution is unreasonable? Explain.
\lines[2]
\end{minipage}

\vspace{3pt}
\textbf{2. A product denominator (Practice 24).} Solve the equation. \quad $\dfrac{x}{x-3}+\dfrac{2x}{x+3}=\dfrac{18}{(x+3)(x-3)}$

Domain: $x\neq$\,\blank[1.2cm] and \blank[1.2cm] \quad Multiply both sides by: \blank[2.6cm]
\lines[4]
Solutions of the cleared equation: \blank[2.4cm] \quad Extraneous: \blank[1.6cm] \quad Solution: \blank[1.6cm]

\vspace{3pt}
\textbf{3. Two denominators (Practice 32).} Which of the following is the solution of $\dfrac{3}{x+1}=\dfrac{2}{x-3}$?

\quad A. $x=-11$ \qquad B. $x=-\tfrac75$ \qquad C. $x=\tfrac75$ \qquad D. $x=11$

Multiply both sides by \emph{both} denominators: \blank[2.6cm]
\lines[3]
Solution: $x=$\,\blank[2cm] \quad Check: \checkmark\ or \texttimes

\newpage
\wspart{IXL}{Practice for the quiz}
Go to IXL, Algebra 2. Type the code in the search box. Do the starred skills first.

\vspace{2pt}
\begin{tabularx}{\textwidth}{l Y Y}
\toprule
\textbf{Code} & \textbf{Skill} & \textbf{Feeds} \\
\midrule
\ixlrow{CHP $\bigstar$}{Solve rational equations}{Quiz 2 (the rate equation), 3, 4; Quiz 1 (full Savvas quiz only)}
\ixlrow{CJC $\bigstar$}{Solve a quadratic equation by factoring}{Quiz 4: the quadratic left after clearing the fractions (Example 4B)}
\ixlrow{7JJ}{Rational functions: asymptotes and excluded values}{Warm-up: the excluded values are the numbers that can never be solutions}
\ixlrow{FEX}{Add and subtract rational expressions}{Warm-up from 4-4: finding a common denominator}
\bottomrule
\end{tabularx}

\textbf{SmartScore 70 or above on each starred ($\bigstar$) skill replaces your quiz score.}

% R3: DOK 3 tied to kept Example 4 (Practice & Problem Solving 13, Higher Order Thinking).
\wspart{Exit question}{Practice 13 (DOK 3)}
\fbox{\parbox{\dimexpr\textwidth-2\fboxsep-2\fboxrule}{%
\textbf{Higher Order Thinking.} Write a rational equation that cannot have 2 or $-6$ as solutions.

\vspace{6pt}
A denominator that is 0 when $x=2$: \blank[2.4cm] \qquad and when $x=-6$: \blank[2.4cm]

\vspace{6pt}
My equation:\lines[2]

\vspace{4pt}
Why can neither 2 nor $-6$ be a solution of your equation? Use the words \emph{domain} and \emph{denominator}.
\lines[3]\vspace{2pt}}}

\vspace{6pt}
\textbf{Before the quiz, finish these sentences.}

To solve a rational equation, I first find the \blank[2.4cm], then I multiply both sides by the \blank[3.6cm].

After I solve, I check each answer. An answer that makes a denominator 0 is an \blank[3cm], so I cross it out.

When two people work together, I \blank[2cm] their \blank[1.6cm], and the time together is \blank[2.4cm] than either time alone.

\end{document}
